\documentclass[runningheads]{llncs}

% ---------------------------------------------------------------
% Include basic ECCV package
% TODO REVIEW: Insert your submission number below by replacing '*****'
% TODO FINAL: Comment out the following line for the camera-ready version
\usepackage[review,year=2026,ID=*****]{eccv}
% TODO FINAL: Un-comment the following line for the camera-ready version
%\usepackage{eccv}

% ---------------------------------------------------------------
% Other packages
\usepackage{eccvabbrv}
\usepackage{graphicx}
\usepackage{booktabs}
\usepackage{amsmath}
\usepackage{multirow}
\usepackage{xcolor}
\usepackage{siunitx}
\sisetup{detect-weight=true,detect-family=true,group-separator={,}}
% axessibility.sty (PDF accessibility for math) relies on pdfTeX-only primitives
% and does not compile under XeTeX/LuaTeX. It is loaded for the pdfLaTeX
% camera-ready build (the template's default engine) and skipped otherwise, with
% no effect on layout or content. To force pdfLaTeX: ``pdflatex main''.
\usepackage{iftex}
\ifPDFTeX
  \usepackage[accsupp]{axessibility}  % Improves PDF readability for those with disabilities.
\fi

% ---------------------------------------------------------------
% Hyperref
% TODO FINAL: Comment out the following line for the camera-ready version
%\usepackage[pagebackref,breaklinks,colorlinks,citecolor=eccvblue]{hyperref}
% TODO FINAL: Un-comment the following line for the camera-ready version
\usepackage{hyperref}
\usepackage{orcidlink}

% A light convenience macro for in-draft notes. Remove before submission.
\newcommand{\todo}[1]{\textcolor{red}{[TODO: #1]}}

\begin{document}

% ---------------------------------------------------------------
\title{Single View Might Be Enough: A Task-Dependent Study of Multi-Perspective Necessity in Geosensing}

\titlerunning{Single View Might Be Enough}

% In review mode, the eccv package auto-replaces the author/institute/running
% heads with anonymous placeholders, so the real names below are only used for
% the camera-ready build. (See eccv.sty review-mode \maketitle override.)
% TODO FINAL: add ORCIDs and correct affiliations for camera-ready.
\author{Ezel Bayraktar\inst{1} \and
Alperen Demirci\inst{1} \and
Erkut Erdem\inst{1} \and
Aykut Erdem\inst{1} \and
Mukhamediyar Amanzhanov\inst{1} \and
Adam Sattout\inst{1}}
\authorrunning{E.~Bayraktar et al.}
\institute{Hacettepe University, Ankara, Turkey}

\maketitle

\begin{abstract}
Pairing overhead (satellite) and street-level imagery is widely assumed to help
vision-language models reason about cities. We test this assumption directly. We
build a multi-city visual question answering benchmark of \num{40} cities in
which each location carries a satellite tile and four street-level views, and
ask four-way multiple-choice questions about eight urban attributes derived from
OpenStreetMap. Using a controlled image ablation---no image, satellite only,
street only, and both---we measure \emph{fusion synergy}, the accuracy a model
gains from both views over its better single view. Across a \num{9}B model we
fine-tune on the data, four off-the-shelf remote-sensing VLMs, two city splits,
and three prompting strategies including extended chain-of-thought, the gain is
marginal: the two-view condition exceeds the better single view by only about
two to three points on average, an order of magnitude below the \num{40}--\num{70}
points the \emph{same} models gain on cross-view questions that require both
images by construction. A paired item-level analysis sharpens the picture: the
two-view model trails a per-item single-view oracle by \num{8}--\num{9} points,
and supplying the second view corrects a wrong answer on roughly \num{1}\% of
items while flipping a correct answer to wrong on roughly \num{10}\%. For these
urban-attribute tasks at commonly served resolutions, the second view contributes
far less than its acquisition cost, and multi-perspective necessity is
task-dependent rather than universal.
  \keywords{Geosensing \and Vision-language models \and Multi-view fusion}
\end{abstract}

% =====================================================================
\section{Introduction}
\label{sec:intro}

A recurring premise in geospatial machine learning is that more perspectives are
better. Overhead imagery captures layout, footprints, and connectivity; ground
imagery captures facades, surfaces, signage, and street furniture. Pairing the
two is intuitively complementary, and a growing line of work collects, aligns,
and jointly models satellite and street-level views for urban understanding.
The implicit expectation is that a model with access to both views should
outperform one restricted to either.

We ask whether this expectation actually holds for concrete urban-attribute
recognition, and find that---for the tasks we study---it largely does not. We
construct a multi-city visual question answering (VQA) benchmark in which every
location is represented by a marked satellite tile and four cardinal
street-level views, and pose four-way multiple-choice questions about eight
urban attributes (land use, building height, urban density, road type, road
surface, junction type, amenity richness, and transit density) grounded in
OpenStreetMap (OSM) metadata~\cite{osm}. Crucially, the same imagery also
supports a family of \emph{cross-view} questions that are unanswerable from a
single view by construction (e.g.\ ``does this street-level photo belong to this
satellite tile?''). These cross-view tasks serve as a built-in positive control:
any sound measurement of view complementarity must register a large benefit
there.

Our central instrument is a controlled image ablation. For each question we
evaluate a model under four input conditions---\textsf{blind} (no image),
\textsf{sat} (satellite only), \textsf{sv} (street only), and \textsf{full}
(both)---and define \emph{fusion synergy} as
$\textsf{full} - \max(\textsf{sat},\textsf{sv})$: the accuracy gained from having
both views over the better single view. Positive synergy is the signature of
genuine cross-view fusion. We complement synergy with \emph{vision contribution},
$\textsf{full} - \textsf{blind}$, which isolates how much the imagery
contributes at all over a text-only prior.

The result is consistent across conditions and, we argue, robustly measured.
On the eight attribute tasks the second view's contribution is marginal: across a
\num{9}B-parameter model we fine-tune on the benchmark, four off-the-shelf
remote-sensing VLMs evaluated zero-shot, two distinct city splits (random
per-city and held-out unseen cities), and three prompting strategies (plain
decoding, a \num{16}k-token explicit reasoning budget, and a four-turn
observe--reason--self-check--commit protocol), having both views improves on the
better single view by only about two to three points on average. A paired
item-level analysis (\cref{sec:itemlevel}) makes the size of this effect concrete:
the two-view model trails a per-item single-view \emph{oracle}---one that simply
picks the better of the two single-view answers for each question---by
\num{8}--\num{9} points, and the joint views contain essentially no information
beyond their better single view. The mechanism is visible in the per-item flips:
adding the second view rescues a wrong answer on about \num{1}\% of items but
overturns a correct one on about \num{10}\%, so the model appears unable to
adjudicate when the two views disagree. Vision itself clearly
matters---\textsf{full}\,$-$\,\textsf{blind} is roughly $+\num{10}$ points---so
the model \emph{is} reading the images; it simply derives a single view's worth of
signal from them. And on the cross-view control tasks the very same models gain
\num{40}--\num{70} points from the second view, demonstrating that the ablation
detects fusion when fusion is genuinely required.

We frame this as a \emph{task-dependent} observation, not a claim that overhead
imagery is uninformative or that multi-view fusion is futile in general. It is a
measured property of single-view-decidable attribute recognition at the satellite
resolutions commonly served in practice. We make this scope explicit throughout,
report bootstrap confidence intervals on every synergy estimate, audit the
benchmark for text-only leakage, and characterize the items where the second view
changes the answer.

Our contributions are:
\begin{enumerate}
  \item A multi-city satellite\,$+$\,street VQA benchmark (\num{40} cities, eight
  OSM-grounded urban-attribute tasks plus cross-view control tasks) with a clean
  location-level held-out split and a controlled four-way image ablation.
  \item A measurement of multi-view fusion synergy that is near zero for all
  eight attribute tasks and robust across model scale, training regime, city
  split, and prompting strategy---with the cross-view tasks as a positive control
  that proves the measurement is sensitive.
  \item An item-level analysis (per-item oracle ceilings, paired significance and
  equivalence tests, and a characterization of the view-flip set) that bounds the
  \emph{latent} complementarity available to any fusion method, reframing ``no
  method helped'' as ``there was little to capture.''
  \item Practical guidance: for these attribute families, serving a single
  (typically cheaper and better-covered) view costs essentially no accuracy, and
  a reusable ablation protocol for deciding when a second view is warranted.
\end{enumerate}

% =====================================================================
\section{Related Work}
\label{sec:related}

\subsubsection{Remote-sensing vision-language models are single-view.}
A wave of instruction-tuned VLMs target remote sensing, including
GeoChat~\cite{geochat2024}, SkySenseGPT~\cite{skysensegpt2024},
LHRS-Bot and LHRS-Bot-Nova~\cite{lhrsbot2024,lhrsbotnova2025}, EarthGPT~\cite{zhang2024earthgpt},
and the multi-sensor generalist EarthDial~\cite{earthdial2025}. These models are
overhead-only: each consumes a single aerial or satellite image, and none ingests
ground-level imagery (\cref{tab:models}). Where ``multi-modal'' appears in this
line of work it refers to multiple \emph{overhead} sensors---optical, SAR, and
infrared, as in EarthGPT---rather than to pairing satellite with street-level
views. The dominant evaluation suites follow suit: RSVQA~\cite{lobry2020rsvqa},
VRSBench~\cite{li2024vrsbench}, and GeoChat's own benchmark all pose questions
over a single overhead image. None isolates the marginal value of a second,
ground-level view, which is the question we study. Architecturally these models
build on LLaVA-style designs~\cite{llava_2023,llava15_2024} with CLIP or SigLIP
encoders~\cite{clip_2021,siglip_2023}; we use four as zero-shot baselines and
discuss their training-data provenance in \cref{sec:setup}.

\subsubsection{Multi-view urban benchmarks.}
A smaller set of benchmarks does pair overhead and ground views for urban
understanding. UrBench~\cite{urbench2025} evaluates multimodal LLMs on
coordinate-aligned street and satellite pairs across fourteen task types, and
reports that models perform \emph{worse} on its cross-view tasks than on either
single view---a pattern its authors attribute to a cross-view reasoning deficit,
and which we revisit here for the narrower case of attribute recognition.
UrbanLLaVA~\cite{urbanllava2025} trains a multi-modal model over satellite,
street-view, OSM, and trajectory inputs; notably its attribute-style tasks (land
use, density) are posed from satellite alone, while its multi-view supervision
targets matching and localization. CityCube~\cite{citycube2026} benchmarks
cross-view \emph{spatial} reasoning (orientation, perspective taking) rather than
attribute recognition. Closest to our finding, the concurrent
CityLens~\cite{citylens2026} benchmarks LVLMs on socioeconomic indicators from
satellite and street-view and reports that street-view alone performs comparably
to using both views and well above satellite alone---a redundancy observation in
the same domain, which our controlled ablation, item-level oracle analysis, and
positive control turn into a quantified and mechanistic account. Our contribution
relative to this group is not a first paired-view dataset, but a controlled
measurement of \emph{when} the second view is needed, with a built-in sensitivity
check.

\subsubsection{Cross-view geo-localization: where both views are necessary.}
Pairing satellite and ground imagery is a well-established setting in cross-view
geo-localization, where the goal is to match a ground photo to an overhead
reference~\cite{durgam2024crossview,zhai2017predicting,zhu2021vigor}. There the
two views are jointly \emph{necessary} by definition---one cannot localize the
ground image without the overhead one---and methods succeed precisely by modeling
the distinct-but-linkable geometric layout the two views
share~\cite{zhang2023geodtr}. This is the regime our cross-view control tasks
occupy, and the contrast with it is exactly our point: complementarity that is
constitutive of localization is largely absent from single-view-decidable
attribute recognition.

\subsubsection{The marginal value of a second modality.}
That adding a modality need not help is known beyond geosensing. Wang
\etal~\cite{wang2020multimodal} show that under naive joint training the best
single-modal network can outperform the multi-modal one, and Liang
\etal~\cite{liang2023quantifying} formalize multimodal interactions into
redundancy, uniqueness, and synergy, demonstrating that low-synergy,
strong-unimodal regimes are real and measurable; our ``fusion synergy'' is the
behavioral analogue of their synergy term. Against this, theory shows multimodal
learning can provably lower achievable risk~\cite{huang2021makes}, and several
aerial$+$ground fusion systems report gains for recognition---a unified near/remote
model for land use and building function~\cite{workman2017unified}, multi-view
scene-classification datasets~\cite{machado2021airound}, and, most directly
counter to our setting, Suel \etal~\cite{suel2021multimodal}, who find that fusing
satellite and street-level imagery measurably improves regression of income and
environmental deprivation. We engage this last result directly in
\cref{sec:discussion}: it differs from ours in task type (continuous regression
versus four-way recognition), label source, and model (a trained CNN versus
zero-shot and fine-tuned VLMs), and our finding refines rather than contradicts
the broad ``fusion helps'' literature by identifying a task regime in which it
does not.

% =====================================================================
\section{The Benchmark}
\label{sec:dataset}

\subsubsection{Construction.}
The benchmark pairs overhead and street-level imagery with multiple-choice
questions about urban attributes, generated from OpenStreetMap (OSM)
metadata~\cite{osm}. Each \emph{location} is sampled within one of \num{40}
cities across \num{32} countries on six continents and is associated with (i) one
satellite tile, (ii) four street-level views at fixed bearings (forward,
backward, left, right), and (iii) OSM-derived attribute labels. Satellite imagery
is sourced per-location by geography---sub-meter aerial orthophotos: ESRI World
Imagery ($\sim$\SIrange{0.3}{0.5}{m/px} in cities) as the global default, NAIP
($\sim$\SI{0.6}{m/px}) in the US, and IGN ($\sim$\SI{0.2}{m/px}) in France, with a
coarse Sentinel-2 fallback ($\sim$\SI{10}{m/px}) reserved for locations that no
high-resolution source covers. In the benchmark every overhead tile resolves to a
sub-meter source (\num{4775} ESRI, \num{305} NAIP, \num{160} IGN; the Sentinel-2
fallback was never triggered), so the redundancy we report below is measured
against genuinely high-resolution overhead imagery rather than a coarse proxy.
The satellite tile carries a marker indicating the queried location; the marker
is part of the input condition and is held fixed across all ablation modes (we
examine its effect in \cref{sec:robustness}).

\subsubsection{Tasks.}
We study eight \emph{attribute} tasks---\textsf{land\_use},
\textsf{building\_height}, \textsf{urban\_density}, \textsf{road\_type},
\textsf{road\_surface}, \textsf{junction\_type}, \textsf{amenity\_richness}, and
\textsf{transit\_density}---each posed as a four-way multiple-choice question
with a single correct letter. These are the questions for which ``do we need both
views?'' is well-posed: each is, in principle, answerable from imagery of a
single location. The dataset additionally contains \emph{cross-view} tasks
(view--view matching at two difficulty levels, in binary and multiple-choice
form, and a camera-direction task) whose answer lives in the \emph{relationship}
between the two views; these are unanswerable from a single view by construction
and we use them only as a positive control (\cref{sec:control}).

\subsubsection{Splits.}
All splits operate at the location level so that no location's questions span
two splits. We hold out \num{10}\% of locations \emph{per city} as a benchmark
test set; this set is identical across training regimes and is disjoint from
training data (verified: zero question-id overlap). For training we use two
strategies that share this benchmark: a \emph{per-city} split (a further fraction
of locations per city to validation) and a \emph{seen/unseen} split (eight cities
held out entirely to validation, to probe cross-city generalization). Question
generation uses a fixed seed for option shuffling, and we keep a single question
per (location, task) pair to avoid location-level inflation. Locations with fewer
than four street-level views are discarded.

\subsubsection{Leakage audit.}
Because the questions are templated, we audit how much of each task is solvable
from text alone by evaluating every model in \textsf{blind} mode (no image; see
\cref{sec:results}). We report \textsf{blind} accuracy per task as a first-class
diagnostic and treat $\textsf{full}-\textsf{blind}$ as the genuinely
vision-attributable signal. During this audit we identified one additional
attribute task (green-space presence) that was answerable from the option text
alone---a model reached \num{100}\% accuracy on it with no image, because the
correct option wording was systematically identifiable independent of its A/B/C/D
position. We \emph{exclude that task entirely} from all results below; the study
reports the eight attribute tasks that survive the audit. We note that across the
retained tasks the correct-option letter is near-uniformly distributed
(A/B/C/D $=$ \num{743}/\num{766}/\num{703}/\num{720} over the \num{2932}
eight-task benchmark questions), so there is no exploitable positional
regularity, and the correct option is not systematically the longest or most
detailed (\cref{sec:robustness}).

\subsubsection{Statistics.}
\Cref{tab:dataset} summarizes the benchmark. The eight-task test set comprises
\num{2932} questions over \num{421} locations across \num{40} cities; each
location supplies one satellite tile and four street-level views.

\begin{table}[t]
  \centering
  \caption{Benchmark test set (eight attribute tasks, held-out \num{10}\% of
  locations per city). Each question has four options and a single correct
  letter; image mode is a single marked satellite tile plus four street-level
  views.}
  \label{tab:dataset}
  \begin{tabular}{@{}lr@{}}
    \toprule
    Property & Value \\
    \midrule
    Cities & 40 \\
    Locations (test) & 421 \\
    Questions (test) & 3135 \\
    Attribute tasks & 8 \\
    Options per question & 4 \\
    Satellite views / location & 1 \\
    Street-level views / location & 4 \\
    \midrule
    \multicolumn{2}{@{}l}{\emph{Questions per task}}\\
    \quad land\_use / urban\_density / road\_type & 421 each \\
    \quad amenity\_richness / transit\_density & 421 each \\
    \quad junction\_type & 334 \\
    \quad road\_surface & 303 \\
    \quad building\_height & 190 \\
    \bottomrule
  \end{tabular}
\end{table}

% =====================================================================
\section{Experimental Setup}
\label{sec:setup}

\subsubsection{Models.}
Our primary model is \textbf{Qwen3.5-9B} fine-tuned on the benchmark's training
split with LoRA~\cite{lora_2022} (rank \num{16}, $\alpha=\num{16}$, applied to
all linear layers including the vision tower) and evaluated on the held-out
benchmark; we train one model per split (per-city and seen/unseen) and report
both. This is the cleanest setting for the view question because the model has
been explicitly optimized on the task distribution---so a failure to use both
views cannot be attributed to a lack of task alignment. To test that the finding
is not specific to our model, we additionally evaluate four \emph{off-the-shelf}
remote-sensing VLMs zero-shot on the overhead half of the benchmark:
GeoChat-7B~\cite{geochat2024}, SkySenseGPT-7B~\cite{skysensegpt2024},
LHRS-Bot-Nova~\cite{lhrsbotnova2025}, and
EarthDial-4B (RGB)~\cite{earthdial2025}. \Cref{tab:models} lists their
architectures and modalities; all four are single-image, overhead-only models.

\begin{table}[t]
  \centering
  \caption{Off-the-shelf remote-sensing VLMs evaluated zero-shot on the overhead
  (satellite-only) half of the benchmark. All four are single-image and
  \emph{overhead-only}; none ingests ground-level imagery. ``OSM data'' marks
  whether the model's training corpus is derived from OpenStreetMap (relevant to
  the contamination caveat in \cref{sec:robustness}).}
  \label{tab:models}
  \setlength{\tabcolsep}{4pt}
  \begin{tabular}{@{}lllll@{}}
    \toprule
    Model & Base LLM & Vision encoder & Input res. & OSM data \\
    \midrule
    GeoChat-7B~\cite{geochat2024}       & Vicuna-1.5-7B & CLIP ViT-L/14 & 504 & No \\
    SkySenseGPT-7B~\cite{skysensegpt2024} & Vicuna-1.5-7B & CLIP ViT-L/14 & 336 & No \\
    LHRS-Bot-Nova~\cite{lhrsbotnova2025}  & Llama-3-8B    & SigLIP-L/14   & 336 & \textbf{Yes} \\
    EarthDial-4B-RGB~\cite{earthdial2025} & Phi-3-Mini    & InternViT-300M & tiled & \textbf{Yes} \\
    \bottomrule
  \end{tabular}
\end{table}

\subsubsection{Image-ablation protocol.}
Every question is evaluated under four input conditions: \textsf{blind} (prompt
only), \textsf{sat} (the marked satellite tile only), \textsf{sv} (the four
street-level views only), and \textsf{full} (satellite tile $+$ four
street-level views). The text of the question and options is byte-identical
across conditions; only the set of supplied images changes. We decode greedily
and parse a single answer letter. For the trained model we use the same image
preprocessing as training (max edge \num{512}\,px).

\subsubsection{Metrics.}
We report per-task accuracy in each condition and two derived quantities:
\emph{fusion synergy} $\;=\textsf{full}-\max(\textsf{sat},\textsf{sv})$ and
\emph{vision contribution} $\;=\textsf{full}-\textsf{blind}$. Synergy
$>0$ indicates the second view adds accuracy a single view cannot; synergy
$\approx 0$ indicates redundancy. \todo{Add bootstrap \num{95}\% CIs and TOST
equivalence margin once the per-item re-analysis is run (Sec.~\ref{sec:results}).}

\subsubsection{Unified harness.}
All models---trained and zero-shot---are routed through one evaluation harness
with a shared multiple-choice prompt, shared image rendering, and shared
answer-letter parsing, so the question text and images are identical across
models and only the forward pass differs. Runs use a single
RTX\,PRO\,6000 (\SI{96}{GB}) GPU.

% =====================================================================
\section{Results}
\label{sec:results}

\subsection{The second view adds little over the better single view}
\label{sec:main}

\Cref{tab:fusion-main} reports the trained Qwen3.5-9B on the held-out benchmark,
per task, under all four conditions, for both splits. Two patterns stand out.
First, \textsf{full} tracks the \emph{better single view} closely: per-task fusion
synergy ranges from $-1.0$ to $+4.5$ points and averages $+\num{1.2}$/$+\num{1.3}$
points (per-city / seen-unseen) across the eight tasks. The gain is positive but
small, and street-level imagery alone (\textsf{sv}) already matches or exceeds the
two-view \textsf{full} condition on several tasks (\textsf{land\_use},
\textsf{road\_type}, \textsf{road\_surface}). Whether this two-to-three-point
margin reflects genuine fusion or merely the model defaulting to the stronger view
is the question we resolve at the item level in \cref{sec:itemlevel}. Second,
vision nonetheless matters: \textsf{full}\,$-$\,\textsf{blind} is large and
positive (e.g.\ $+\num{22}$ on \textsf{transit\_density}, $+\num{32}$ on
\textsf{amenity\_richness}), so the model reads the imagery---it simply extracts a
single view's worth of information from it.

\begin{table}[t]
  \centering
  \caption{Trained Qwen3.5-9B on the held-out benchmark, per attribute task and
  input condition (accuracy, \%). \textbf{Syn.} $=\textsf{full}-\max(\textsf{sat},
  \textsf{sv})$ is fusion synergy. Synergy is small on every task under both
  splits; \textsf{full} stays close to the better single view.}
  \label{tab:fusion-main}
  \setlength{\tabcolsep}{4.5pt}
  \begin{tabular}{@{}lcccc@{\;}c@{\qquad}cccc@{\;}c@{}}
    \toprule
    & \multicolumn{5}{c}{\textbf{per-city split}} & \multicolumn{5}{c}{\textbf{seen/unseen split}} \\
    \cmidrule(lr){2-6}\cmidrule(lr){7-11}
    Task & blind & sat & sv & full & Syn. & blind & sat & sv & full & Syn. \\
    \midrule
    land\_use         & 68.6 & 75.5 & 75.1 & 75.3 & $-0.2$ & 63.4 & 72.0 & 72.4 & 73.9 & $+1.4$ \\
    building\_height  & 45.3 & 67.4 & 66.3 & 66.8 & $-0.5$ & 50.5 & 63.7 & 66.3 & 66.8 & $+0.5$ \\
    urban\_density    & 54.6 & 61.8 & 62.7 & 65.1 & $+2.4$ & 53.2 & 71.5 & 64.8 & 74.6 & $+3.1$ \\
    junction\_type    & 61.4 & 73.4 & 68.6 & 75.1 & $+1.8$ & 62.0 & 74.0 & 71.3 & 76.3 & $+2.4$ \\
    amenity\_richness & 26.4 & 53.4 & 56.5 & 58.7 & $+2.1$ & 25.7 & 53.0 & 55.6 & 55.6 & $+0.0$ \\
    road\_type        & 68.6 & 72.4 & 77.4 & 76.5 & $-1.0$ & 68.6 & 69.6 & 76.7 & 76.7 & $+0.0$ \\
    road\_surface     & 95.7 & 95.0 & 93.4 & 95.7 & $+0.7$ & 95.7 & 95.0 & 95.4 & 95.0 & $-0.3$ \\
    transit\_density  & 28.3 & 43.7 & 45.6 & 50.1 & $+4.5$ & 26.4 & 45.8 & 44.9 & 49.4 & $+3.6$ \\
    \midrule
    \textbf{Mean}     & 56.1 & 67.8 & 68.2 & 70.4 & $\mathbf{+1.2}$ & 55.7 & 68.1 & 68.4 & 71.1 & $\mathbf{+1.3}$ \\
    \bottomrule
  \end{tabular}
\end{table}

\subsection{The same models fuse when fusion is required (positive control)}
\label{sec:control}

A reader should ask whether our ablation could detect fusion at all---perhaps the
small attribute-task synergy reflects a measurement that is simply insensitive to
the second view. It does not. \Cref{tab:control} reports the identical model and
harness on the cross-view control tasks, whose answers require comparing the two
views. Here synergy is large---$+\num{17}$ to $+\num{70}$ points---because a single
view is, by construction, insufficient. The contrast with \cref{tab:fusion-main}
is the crux of the paper: the same instrument that registers a two-to-three-point
gain on the attribute tasks registers a forty-to-seventy-point gain the moment the
task genuinely requires both views. The marginal synergy on attributes is
therefore a property of those tasks, not an artifact of an insensitive test.

\begin{table}[t]
  \centering
  \caption{Positive control: the same trained Qwen3.5-9B (per-city) on cross-view
  tasks that require both views by construction. Synergy is large and positive,
  confirming the ablation detects fusion when it is needed.}
  \label{tab:control}
  \setlength{\tabcolsep}{5pt}
  \begin{tabular}{@{}lccccc@{}}
    \toprule
    Cross-view task & blind & sat & sv & full & Syn. \\
    \midrule
    camera\_direction      & 25.2 & 24.9 & --   & 42.3 & $+17.3$ \\
    mismatch\_binary\_easy & 46.6 & 51.1 & 52.7 & 97.4 & $+44.7$ \\
    mismatch\_binary\_hard & 55.8 & 45.8 & 44.4 & 90.7 & $+44.9$ \\
    mismatch\_mcq\_easy    & 20.7 & 26.4 & 23.5 & 96.2 & $+69.8$ \\
    mismatch\_mcq\_hard    & 22.3 & 23.3 & 22.8 & 86.0 & $+62.7$ \\
    \bottomrule
  \end{tabular}
\end{table}

That the cross-view tasks genuinely encode cross-view information---rather than
being solved by some shortcut our VLM happens to exploit---is corroborated by an
independent, purpose-built model. We evaluate Sample4Geo~\cite{deuser2023sample4geo},
a Siamese ConvNeXt trained for street$\leftrightarrow$satellite retrieval, on the
mismatch family by embedding the satellite tile and each street-view set and
scoring by cosine similarity. This is a different model class and metric from our
MCQ VQA and the numbers are not comparable to \cref{tab:control}; we report it only
to confirm the tasks' nature. Transferred zero-shot (its VIGOR checkpoint, the
closest urban domain), it reaches \num{0.65} accuracy on the binary
same-place/different-place questions---well above the \num{0.50} chance rate---so a
model that only measures cross-view similarity already carries real signal on these
items. It is near chance on the four-way variants (\num{0.27}--\num{0.31} vs.\
\num{0.25}), reflecting how much harder it is to single out the correct nearby
scene from three plausible distractors across the viewpoint gap. The attribute
tasks, by contrast, offer no such cross-view signal to recover.

\subsection{The finding holds zero-shot across off-the-shelf RS-VLMs}
\label{sec:zeroshot}

To rule out that the redundancy is an artifact of our particular model, we
evaluate four off-the-shelf remote-sensing VLMs zero-shot on the overhead
(satellite-only) half of the benchmark (\cref{tab:zeroshot}). None was trained on
our task, our four-way MCQ format, or our marker overlay. They span CLIP- and
SigLIP-based encoders, three LLM families, and single-crop versus tiled
preprocessing. The trained Qwen3.5-9B in its satellite-only mode is the
like-for-like single-overhead-view comparison; it outperforms the best zero-shot
model (EarthDial, $\num{39.9}\%$) by $+\num{26.7}$ points overall, as expected
for a model fine-tuned on the distribution. The four zero-shot models are
satellite-only and therefore inform the single-view baseline and the
contamination discussion (\cref{sec:robustness}), not the synergy claim itself,
which is a within-model comparison.

\begin{table}[t]
  \centering
  \caption{Off-the-shelf RS-VLMs, zero-shot, single marked satellite tile,
  per-task accuracy (\%) over the eight-task benchmark ($n=\num{2932}$). Last
  column: our trained Qwen3.5-9B in satellite-only mode (per-city), the
  like-for-like single-overhead-view reference. LHRS-Bot-Nova and EarthDial are
  trained on OSM-derived text and their scores should be read as an
  \emph{upper bound} on visual skill (see \cref{sec:robustness}).}
  \label{tab:zeroshot}
  \setlength{\tabcolsep}{4pt}
  \begin{tabular}{@{}lccccc@{}}
    \toprule
    Task & GeoChat & SkySenseGPT & LHRS-Nova & EarthDial & \textit{Ours (sat)} \\
    \midrule
    land\_use         & 48.7 & 53.4 & 46.3 & 47.5 & \textit{75.5} \\
    building\_height  & 30.5 & 34.7 & 28.9 & 32.6 & \textit{67.4} \\
    urban\_density    & 13.5 & 14.5 & 17.8 & 30.4 & \textit{61.8} \\
    road\_type        & 10.7 & 18.1 & 26.1 & 46.1 & \textit{72.4} \\
    road\_surface     & 75.6 & 80.2 & 58.7 & 58.4 & \textit{95.0} \\
    junction\_type    &  8.4 & 30.8 & 35.3 & 52.7 & \textit{73.4} \\
    amenity\_richness & 13.5 & 12.6 & 17.8 & 25.4 & \textit{53.4} \\
    transit\_density  & 21.4 & 26.1 & 34.2 & 29.7 & \textit{43.7} \\
    \midrule
    \textbf{Overall}  & 26.2 & 32.0 & 32.4 & \textbf{39.9} & \textit{\textbf{66.6}} \\
    \bottomrule
  \end{tabular}
\end{table}

\subsection{No prompting strategy recovers complementary signal}
\label{sec:strategies}

If the two views \emph{were} complementary but the model simply failed to
combine them, a stronger reasoning procedure should surface the latent signal. It
does not. \Cref{tab:strategies} reports three strategies---plain decoding, a
\num{16}k-token explicit-thinking budget, and a four-turn
\emph{observe$\rightarrow$reason$\rightarrow$self-check$\rightarrow$commit}
conversation---on the eight attribute tasks, here using an AWQ-quantized
Qwen3.5-9B so the long-context reasoning runs fit a single GPU. More deliberation
raises accuracy modestly---the model gets better at each task in every
view---but synergy stays \emph{negative} throughout, and indeed \textsf{sv} alone
is the single strongest condition under every strategy. The extra reasoning
accrues to the single-view conditions just as much as to \textsf{full}, so what
improves is the model's reading of one view, not its ability to combine two. This
mirrors the item-level finding of \cref{sec:itemlevel}: the obstacle is
adjudicating between redundant views, which more thinking does not remove.

\begin{table}[t]
  \centering
  \caption{Prompting strategies on the eight attribute tasks (Qwen3.5-9B-AWQ,
  $n=\num{2932}$, per-task means). \textbf{Syn.} $=\textsf{full}-\max(\textsf{sat},
  \textsf{sv})$; \textbf{Vis.} $=\textsf{full}-\textsf{blind}$. Reasoning improves
  accuracy but not synergy---the second view remains redundant, and street-view
  alone (\textsf{sv}) is the single strongest condition in every strategy.}
  \label{tab:strategies}
  \setlength{\tabcolsep}{6pt}
  \begin{tabular}{@{}lccccrr@{}}
    \toprule
    Strategy & blind & sat & sv & full & Syn. & Vis. \\
    \midrule
    plain                 & 41.0 & 46.5 & 54.2 & 51.7 & $-2.8$ & $+10.8$ \\
    think-16k             & 41.4 & 48.2 & 54.0 & 52.5 & $-2.0$ & $+11.2$ \\
    multistep (4-turn)    & 43.4 & 50.3 & 55.0 & 53.2 & $-2.1$ & $+9.8$ \\
    \bottomrule
  \end{tabular}
\end{table}

\subsection{Item-level analysis: the second view interferes more than it helps}
\label{sec:itemlevel}

A small positive average could still hide genuine fusion on a subset of items, or
could be an artifact of the model defaulting to its stronger view. We resolve this
by joining the predictions of all four conditions at the item level (every
benchmark question carries a stable identifier shared across conditions) and
analyzing the \num{2932} fully paired eight-task items directly.

\subsubsection{A per-item single-view oracle beats the two-view model.}
For each item we define an \emph{oracle} that is allowed to pick the better of the
two single-view answers---a model that always chooses the right view, but never
fuses. \Cref{tab:oracle} compares the two-view \textsf{full} condition against
this oracle. The two-view model trails the per-item single-view oracle by
$\num{9.3}$ points (per-city) and $\num{7.9}$ points (seen/unseen), with paired
bootstrap \num{95}\% confidence intervals that exclude zero by a wide margin
($[\num{-10.4},\num{-8.2}]$ and $[\num{-9.1},\num{-6.8}]$, \num{10000} resamples).
Moreover, the fraction of items solved by \emph{either} single view (the union
ceiling) equals the oracle exactly: there is no pool of items that only the
\emph{combination} of views answers correctly. In other words, the two views
carry essentially no answer-relevant information beyond their better single view,
and a system that simply routed each question to the right view would outperform
one that sees both.

\subsubsection{The second view overturns more correct answers than it rescues.}
The mechanism behind the oracle gap is an asymmetry in how the second view changes
answers. Relative to the better single view per item, adding the second view
corrects a previously wrong answer on only \num{0.9}\% / \num{1.2}\% of items
(per-city / seen-unseen) but flips a previously correct answer to wrong on
\num{10.2}\% / \num{9.1}\%---roughly a ten-to-one ratio against. The model does
not reliably defer to whichever view is right; when the views disagree, the second
view drags down items the better view had already answered. This is consistent
with the prompting results of \cref{sec:strategies}: no amount of explicit
reasoning recovers a complementary signal, because the limiting factor is
adjudication between redundant views, not reasoning depth.

\subsubsection{Aggregate synergy is small and positive, not zero.}
For completeness, \cref{tab:oracle} also reports the pooled fusion synergy with
its bootstrap CI: $+\num{2.2}$ $[\num{0.9},\num{3.3}]$ (per-city) and $+\num{2.8}$
$[\num{1.5},\num{3.7}]$ (seen/unseen). These pooled values are item-weighted and
thus slightly exceed the $+\num{1.2}$/$+\num{1.3}$ unweighted task means of
\cref{tab:fusion-main}, because the larger tasks happen to carry the larger
synergies. We do not claim the gain is zero---the interval excludes zero. We claim
only that it is small: an order of magnitude below the $+\num{40}$ to $+\num{70}$
points the same model gains on the cross-view control tasks (\cref{tab:control}),
and net-negative against optimal single-view selection. Whether a $\sim$\num{2}-point average gain justifies acquiring,
storing, and serving a second imagery modality is the practical question we take
up in \cref{sec:discussion}.

\begin{table}[t]
  \centering
  \caption{Item-level analysis on the \num{2932} paired eight-task questions.
  \textbf{Oracle} picks the better single-view answer per item; \textbf{Union} is
  the fraction either single view answers correctly. \textbf{full$-$oracle} and
  \textbf{synergy} carry paired bootstrap \num{95}\% CIs (\num{10000} resamples).
  \textbf{Help}/\textbf{Hurt} are the percentages of items on which the second
  view flips the better single-view answer right and wrong, respectively. The
  two-view model trails the per-item single-view oracle on both splits.}
  \label{tab:oracle}
  \setlength{\tabcolsep}{5pt}
  \begin{tabular}{@{}lcccccc@{}}
    \toprule
    Split & full & oracle & full$-$oracle & synergy & Help & Hurt \\
    \midrule
    per-city    & 69.5 & 78.8 & $-9.3$ \footnotesize$[-10.4,-8.2]$ & $+2.2$ \footnotesize$[0.9,3.3]$ & 0.9\% & 10.2\% \\
    seen/unseen & 70.3 & 78.2 & $-7.9$ \footnotesize$[-9.1,-6.8]$ & $+2.8$ \footnotesize$[1.5,3.7]$ & 1.2\% & 9.1\% \\
    \bottomrule
  \end{tabular}
\end{table}

% =====================================================================
\section{Robustness and Confounds}
\label{sec:robustness}

\todo{Each item below pairs an objection with the concrete check that addresses
it; checks marked \textsf{[planned]} require limited new inference.}

\subsubsection{Satellite resolution.}
Our claim is bounded to commonly served resolutions, but most cities here are
covered by \emph{sub-meter} aerial orthophotos (NAIP $\sim$\SI{0.6}{m/px}, IGN
$\sim$\SI{0.2}{m/px}, ESRI $\sim$\SIrange{0.3}{0.5}{m/px}); Sentinel-2
($\sim$\SI{10}{m/px}) is only a fallback. The redundancy is thus measured with
genuinely high-resolution overhead imagery for the majority of locations, not a
coarse proxy. For tasks that are structurally ground-level (e.g.\ road surface,
signage-driven amenities), no overhead resolution recovers the cue; for tasks
where higher GSD might plausibly help (building height, urban density) we report
a targeted high-resolution sub-study \textsf{[planned]}.

\subsubsection{The satellite marker.}
The satellite tile carries a location marker. We verify the results are not an
artifact of it via an unmarked-tile ablation \textsf{[planned]} and via the
attention analysis below.

\subsubsection{Single-image vs.\ four-view asymmetry.}
\textsf{sv} supplies four images and \textsf{sat} one, so part of \textsf{sv}'s
strength could be image count rather than viewpoint. We test a single
street-level angle versus the four-view set \textsf{[planned]}; our thesis is
explicitly that the \emph{cheaper-to-serve single view suffices}, so street-level
richness is part of why it does, not a flaw in the comparison.

\subsubsection{Text-only leakage.}
We report \textsf{blind} accuracy per task (\cref{sec:results}) as the empirical
ceiling of what is guessable without images, which subsumes positional, option
length, and prior effects. The one task that was fully text-solvable was removed
(\cref{sec:dataset}); we flag remaining tasks with elevated \textsf{blind}
accuracy (notably road surface) and rank tasks by vision contribution rather than
raw accuracy.

\subsubsection{OSM provenance / contamination.}
Our labels are OSM-derived. Two of the four zero-shot models---LHRS-Bot-Nova and
EarthDial---are trained on OSM-derived text (LHRS-Align captions OSM tags across
\num{9259} cities; EarthDial's pretraining QA is built from OSM-labeled imagery),
so their scores may reflect shared OSM priors and we treat them as an upper bound
on visual skill. GeoChat and SkySenseGPT are \emph{not} OSM-derived. Critically,
the synergy claim is a \emph{within-model difference} (\textsf{full} vs.\ single
view), so any uniform OSM prior cancels in the difference. We also report whether
removing geographic identifiers from the prompt changes \textsf{blind} accuracy
\textsf{[planned]}.

\subsubsection{Interpretability (corroboration only).}
Our view-importance claims rest on \emph{causal} input ablation, not attention.
As corroboration, we measure the cross-modal attention mass the model places on
the satellite token versus the street-level tokens when committing to an answer,
reported per task and per layer with normalization for image-token count (a
single satellite image vs.\ four street-level images), and explicitly framed as
consistent-with rather than evidence-of importance, in line with known limits of
attention as explanation~\cite{jain2019attention,wiegreffe2019attention}.
\todo{Insert the normalized attention-vs-causal-contribution correlation once
computed from the existing attention logs.}

% =====================================================================
\section{Discussion}
\label{sec:discussion}

The picture that emerges is not that overhead imagery is useless, nor that fusion
never helps. It is that complementarity is a property of the \emph{task}. When a
question is answerable from a single location's imagery---what land use is here,
how tall are the buildings, is the road paved---the two views encode largely the
same answer-relevant information, and a model reads it off whichever view is more
convenient. When a question is \emph{about the relationship} between the
views---does this photo belong to this tile, which way is the camera
facing---the second view is indispensable, and our models exploit it readily. The
eight attribute tasks fall in the first regime; the cross-view tasks in the
second.

This has a concrete, practical consequence. Street-level and high-resolution
satellite imagery differ sharply in acquisition cost, coverage, licensing, and
latency. For the attribute families studied here, the second view buys a gain of
only about two to three points on average---and, supplied naively, it degrades
roughly one item in ten that the better single view already answered. A
deployment that serves a single, well-chosen view therefore forgoes very little
accuracy while halving its imagery footprint; if both views are available, the
evidence favors \emph{routing} each question to one view over feeding the model
both. The synergy, oracle-gap, and vision-contribution measurements we use form a
reusable instrument for making that call on a new task before committing to
dual-view collection.

% =====================================================================
\section{Conclusion}
\label{sec:conclusion}

We measured, rather than assumed, the value of a second perspective for urban
attribute recognition. Through a controlled four-way image ablation on a
multi-city satellite$+$street VQA benchmark, we found that across model scale,
training regime, city split, and prompting strategy the second view adds only a
few points over the better single view---and that a two-view model trails a
per-item single-view oracle by \num{8}--\num{9} points, overturning correct
answers about ten times as often as it rescues wrong ones. The same models gain
\num{40}--\num{70} points on cross-view tasks that genuinely require both views,
so the effect is a property of the tasks, not of the instrument. Multi-perspective
necessity is task-dependent; for these urban-attribute tasks, a single
well-chosen view might be enough.

% ---------------------------------------------------------------
% TODO FINAL: remove for camera-ready; required for review.
\clearpage

%\section*{Acknowledgements}
% TODO FINAL: add acknowledgements (no section number).

\bibliographystyle{splncs04}
\bibliography{eollm,main}
\end{document}
